\clearpage
\fancyhead[L]{\small HARDING. HIMMELSATLAS.}
\SourceLeaf{515}
\begin{center}\small Göttingische gelehrte Anzeigen. No. 61. pp. 601..607. 17 April 1809.\\[3pt]\rule{.72\linewidth}{.3pt}\end{center}

\textsc{K. L. Harding}. \emph{New Star Atlas}. \emph{First Instalment}, containing sheets I. II. V and IX. Available from the author and through \textsc{Friedrich Perthes} on commission in Hamburg, 1809.

We are pleased to announce the first fruits of an undertaking that experts and lovers of astronomy have eagerly awaited for several years. The hardships of the times put serious obstacles in its way, but fortunately they did not hold up its progress. The need for highly detailed charts of the sky that were as accurate as possible had often been felt before, especially when comets appeared. Since the discovery of Ceres and the three other new sister planets, this need has become urgent for anyone wishing to observe, or even just to see, these remarkable celestial bodies. Anyone with an interest in astronomy must therefore be delighted to see this gap filled, and filled as well as these first four sheets show. Our distinguished colleague decided on this undertaking more than five years ago, soon after the discovery of Ceres and Pallas. His plan at first was admittedly somewhat more limited and called for charting only those zones of the sky in which the two planets could appear to us. Marking out these zones gave rise at the time to an interesting analytical problem that Professor \textsc{Gauss} solved in the Monatliche Correspondenz (August 1804 [p. 106 of this volume]). Professor \textsc{Harding} began his work under the happiest auspices. The discovery of a third new planet ({\LatitudeSymbols οὐ τύχης ἀλλ' ἀρετῆς εὐτυχούσης ἔργον}, as \textsc{Delambre} says in Plutarch’s words) was, as it were, an advance reward for the extensive work he undertook. Yet this very discovery and Dr \textsc{Olbers}’s successful discovery of a fourth planet in 1807 already made it necessary to extend the plan. Convinced of the importance of the undertaking, and in response to repeated requests from other astronomers, Professor \textsc{Harding} dropped the original restrictions altogether and extended the project to the whole sky of the fixed stars, as far as it is of interest to European observers and as far as adequate means of charting it are available. The only traces left of that first limited plan are therefore in the numbering and arrangement of the charts. It is of course also most convenient for those interested in this new representation of the starry sky that the sheets covering parts of the zodiac of the new planets appear first, so that the most pressing need is met as soon as possible.

Professor \textsc{Harding} has included \emph{all} stars for which \emph{reliable} determinations or observations exist. He has also entered, from his own observations, a considerable number of stars that do not yet appear in any catalogue, and he will soon add these determinations to the atlas. Given how much has been done, especially recently, to improve our knowledge of the fixed stars, readers can judge for themselves how far this new atlas must surpass all its predecessors in the richness of its contents. The constellation Aquarius, for example, has 108 stars in \textsc{Flamsteed}’s charts and 343 in the large charts by \textsc{Bode}, whereas we count 1444 stars on Professor \textsc{Harding}’s sheets IX and VIII. The latter is intended for the second instalment, though we already have an engraved copy in hand. It is hard to decide which is the great\SourceLeaf{516}er merit of these charts, this wealth of material itself or the unprecedented care with which it has been used. All earlier makers of star charts plotted the stars as given in the sources they considered best. Some, like Professor \textsc{Bode} in his atlas, succeeded by careful comparisons, conjectures and perceptive criticism in removing many errors of all kinds, which were inevitable in catalogues of fixed stars containing such an enormous number of figures. Yet many errors that these methods could not detect were bound to go unnoticed and pass into the charts. That is why earlier star charts contain so many stars that are nowhere in the sky and never have been. For example, Professor \textsc{Harding} has found no fewer than 15 such stars among \textsc{Bode}’s 343 stars in Aquarius and proportionately many more in some other constellations. The only way to avoid these errors was to check \emph{every} star carefully \emph{in the sky itself} before entering it on the charts. For an undertaking on this scale, that almost seemed beyond the powers of one person, yet Professor \textsc{Harding} considered it indispensable in making a new atlas. Accordingly, there is not a single star in these charts whose actual presence in its place Professor \textsc{Harding} had not first verified for himself.

These general remarks on the spirit in which the undertaking is being carried out make it easy to judge its scientific value. We would like to add a few words about the appearance of the charts. Professor \textsc{Harding} thought it best to use a much smaller format than \textsc{Bode}’s, partly for ease of use and partly to avoid raising the price too much --- both considerations that matter greatly to astronomers. The sheets are $16\frac{3}{4}$ Paris inches high and $20\frac{1}{2}$ inches wide, so their area is only about a third of that of \textsc{Bode}’s charts. On the first four sheets, where the declinations do not extend beyond 30 degrees north and south, it seemed unnecessary, rightly, to choose a proper projection. Degrees of right ascension and declination are therefore drawn to the same scale throughout (about $5\frac{6}{7}$ Paris lines). This also makes the charts more convenient to use in several ways, far outweighing the slight distortion, which is barely noticeable even at the highest declinations. Each sheet thus covers 42 degrees of right ascension and $34\frac{1}{2}$ degrees of declination. In right ascension the first sheet extends from $359^\circ$ to $41^\circ$, the second from $39^\circ$ to $81$, and so on, so that the first joins up again with the ninth, which runs from $319^\circ$ to $1^\circ$. The grid showing each degree of right ascension and declination has been kept for ease of use. Apart from this grid and the boundaries and names of the constellations, the charts contain only what someone buying them for actual use would look for, as opposed to someone collecting pretty pictures. That is, they contain the stars and the designations generally accepted for them, the letters of \textsc{Bayer} or the numbers of \textsc{Flamsteed}. The figures of the constellations themselves are entirely out of use in the present state of the science and would only have disfigured the background of the charts in the eyes of an expert. They have rightly been left out altogether, which is precisely what makes it possible to show so many stars clearly and neatly. On average one will be able to count more than 2000 stars per sheet. If the whole work consists, as the author plans, of 24 sheets, most of which are already fully drawn, the atlas will contain about 50000 stars.

Some readers may wish that the author had added a grid for the degrees of longitude and latitude as well as the one for right ascension and declination, \SourceLeaf{517}especially to make it easier to plot the positions of planets, which astronomical ephemerides usually give only by longitude and latitude. This grid was deliberately left out to keep the framework of black lines from interfering any more than necessary with the visual impression made by the star patterns in the charts. In its place, however, Professor \textsc{Harding} will give a very convenient, carefully calculated table in the text to be added at the end. With it any amateur can easily draw this grid by hand with a pen in their own copy, perhaps in a different colour if they wish, to serve both purposes at once. Perhaps the author would be willing to make this useful table available earlier in a suitable place for those who own the first instalment.

It remains only for us to name the regions of the sky shown on the four sheets of the first instalment. To save as much space as possible and avoid needless repetition, the charts are divided by degrees of right ascension and declination instead of by constellation. Most constellations are therefore split across the charts. The informed amateur hardly needs to be told that this makes no difference at all. Sheet No. I. extends in declination from $13^\circ$ north to $21\frac{1}{2}^\circ$ south; it thus contains parts of Pisces and Aries, most of Cetus, and a small part of Eridanus. On sheet No. II. from $28\frac{1}{2}^\circ$ north to $6^\circ$ south, one finds most of the rest of Aries and Cetus, almost all of Taurus and George’s Harp, and parts of Eridanus and Orion. Sheet No. V. from $27^\circ$ north to $7\frac{1}{4}^\circ$ south shows a large part of Leo, parts of Leo Minor and Crater, and most of Virgo and Berenice’s Hair. Finally sheet No. IV. from $5^\circ$ north to $29\frac{1}{4}^\circ$ south shows part of Capricorn, the Balloon, the Southern Fish, the Sculptor’s Workshop, Cetus, Pisces, Pegasus and Equuleus, and most of Aquarius. --- We end this notice with the wish that all those who care about the flourishing of the sciences and the encouragement of genuinely German perseverance in truly useful work will now do their part to support and sustain this undertaking, which we cannot recommend highly enough. May it bear abundant fruit for this lofty science in years to come, as the finest reward for its worthy author.
